\section{CQ6}

\begin{fancyquote}
    \textbf{CQ6a)} How many different MQTT clients sent atleast one message to erase a previously retained value in a topic in the public HiveMQ broker?\\
    \textbf{CQ6b)} How many of the clients found in CQ6a) connected to the HiveMQ broker with a client identifier strictly longer than 7 bytes (>7)?
\end{fancyquote}

To solve the question, we first identified the IP addresses of the public HiveMQ broker from DNS resolution:
\begin{lstlisting}
tshark -r B.pcapng -Y 'dns.flags.response == 1 && dns.qry.name contains "hivemq"' -T fields -e dns.qry.name -e dns.a
\end{lstlisting}
\[
35.158.43.69,\; 35.158.34.213,\; 18.192.151.104
\]

Then we considered only MQTT traffic directed to those IP addresses.

In MQTT, a retained value is erased by sending a \texttt{PUBLISH} message with:
\begin{itemize}
    \item \texttt{mqtt.msgtype == 3} (\texttt{PUBLISH})
    \item retain flag set
    \item empty payload
\end{itemize}

Therefore, the solution was structured in two parts:
\begin{enumerate}
    \item map each TCP stream to the corresponding MQTT client identifier using \texttt{CONNECT} packets;
    \item find the TCP streams containing retained \texttt{PUBLISH} messages with empty payload, then associate those streams with their client identifiers.
\end{enumerate}

\subsection*{1. Extract MQTT client identifiers from CONNECT packets}

We extracted all MQTT \texttt{CONNECT} packets directed to the HiveMQ broker and saved the mapping \texttt{tcp.stream $\rightarrow$ mqtt.clientid}:

\begin{lstlisting}
tshark -r B.pcapng -Y '(ip.dst == 35.158.43.69 || ip.dst == 35.158.34.213 || ip.dst == 18.192.151.104) && mqtt.msgtype == 1' -T fields -e tcp.stream -e mqtt.clientid | sort -u > stream_clientid.txt
\end{lstlisting}

This produced the following stream-to-client mapping:

    \begin{verbatim}
    1   pcbtjyj
    3   cpoepjzkhibxgjiu
    8   dzcxnwdqef
    15  tukvxesuhe
    20  fcthvjikxjul
    22  giuxfijwus
    23  pohujuv
    43  atqeblw
    44  rolsmcwhshlbc
    \end{verbatim}


\subsection*{2. Find retained PUBLISH messages with empty payload}

Then we extracted all retained MQTT \texttt{PUBLISH} packets directed to the HiveMQ broker:

\begin{lstlisting}
tshark -r B.pcapng -Y "($BROKER_DST) && mqtt.msgtype == 3 && mqtt.retain == 1" -T fields -e frame.number -e tcp.stream -e mqtt.topic -e mqtt.msg -e mqtt.msg_text > retained_candidates.txt
\end{lstlisting}

In the output, packets used to erase a retained value are those where the payload field is missing, shown by \texttt{<MISSING>}.  
Thus we extracted the corresponding TCP streams:

\begin{lstlisting}
awk -F '\t' '($4 == "<MISSING>" || $4 == "") {print $2}' retained_candidates.txt | sort -u > erase_streams.txt
\end{lstlisting}
The resulting streams were:
    \begin{verbatim}
    1
    22
    23
    43
    44
    \end{verbatim}

\subsection*{3. Associate erase streams with client identifiers}

Finally, we associated the streams found above with the client identifiers extracted from the \texttt{CONNECT} packets:

\begin{lstlisting}
sort erase_streams.txt > erase_streams_sorted.txt
sort stream_clientid.txt > stream_clientid_sorted.txt
join erase_streams_sorted.txt stream_clientid_sorted.txt > cq6_clients.txt
\end{lstlisting}

So the distinct MQTT clients that sent at least one retained-delete message are:
\[
\{\texttt{pcbtjyj},\; \texttt{giuxfijwus},\; \texttt{pohujuv},\; \texttt{atqeblw},\; \texttt{rolsmcwhshlbc}\}
\]

\textbf{\huge Answer CQ6a: 5}

\subsection*{4. Count client identifiers longer than 7 bytes}

Among the 5 clients found in CQ6a:
\begin{itemize}
    \item \texttt{pcbtjyj} has length 7
    \item \texttt{giuxfijwus} has length 10
    \item \texttt{pohujuv} has length 7
    \item \texttt{atqeblw} has length 7
    \item \texttt{rolsmcwhshlbc} has length 13
\end{itemize}

Therefore, the client identifiers strictly longer than 7 bytes are:
\[
\texttt{giuxfijwus},\; \texttt{rolsmcwhshlbc}
\]

\textbf{\huge Answer CQ6b: 2}